\chapter{Introduction}
\label{chapter:introduction}

Multilingual communication is a fundamental requirement in many modern contexts, including international conferences, governmental institutions, and global online events. Simultaneous interpretation plays a critical role in enabling such communication by allowing spoken content in a source language to be rendered into one or more target languages with minimal delay. Traditionally, this task has been performed by highly trained human interpreters supported by specialized audio infrastructure \cite{Seeber2015CognitiveLoad}. While this approach delivers high-quality results, it is costly, difficult to scale, and constrained by logistical and availability limitations \cite{dhawan2022speech2speechchallenges, gencheva2023features}.

Recent advances in artificial intelligence have opened new opportunities for automating parts of the interpretation process. Progress in automatic speech recognition, neural machine translation, and text-to-speech synthesis has made it technically feasible to construct end-to-end systems capable of transforming spoken language into translated speech in near real time \cite{popescu-belis2025s2streview}. These developments have sparked growing interest in AI-based simultaneous interpretation systems that aim to reduce operational costs, improve accessibility, and support multilingual communication at scale \cite{ozkaya2023transformative}.

Despite this progress, real-time AI-based interpretation remains a challenging engineering problem. Unlike offline translation or post-event transcription, simultaneous interpretation imposes strict latency requirements that directly affect usability \cite{zheng-etal-2020-fluent}. Spoken input arrives continuously, intermediate representations are inherently unstable, and delays introduced at each processing stage can accumulate rapidly. Furthermore, interpretation systems must operate reliably over extended sessions, support multiple target languages, and integrate seamlessly with audio input and output mechanisms.

Most existing AI-based interpretation approaches can be broadly categorized into either end-to-end speech-to-speech models \cite{jia2022translatotron2} or cascaded pipelines composed of separate speech recognition, translation, and synthesis components \cite{popescu-belis2025s2streview}. While end-to-end models offer conceptual simplicity, they often lack the flexibility and controllability required for real-time deployment. Cascaded architectures, on the other hand, enable modularity and practical integration but introduce challenges related to segmentation, synchronization, and stability across components. As a result, many current systems struggle to balance low latency with coherent and reliable output in live settings \cite{sperber2020endtoendpromise, Sethiya2023EndtoEndSurvey, Etchegoyhen2022CascadeDirectCase}.

In addition to algorithmic challenges, practical deployment considerations further complicate the design of real-time interpretation systems. Professional interpretation environments rely on standardized audio formats, routing mechanisms, and dedicated hardware infrastructure \cite{conference_interpreting_new_tech}. However, many AI-based solutions are developed primarily for software-only or consumer-oriented use cases \cite{dhawan2022speech2speechchallenges} , limiting their applicability in professional contexts \cite{ozkaya2023transformative}. Moreover, existing platforms are often provider-locked, restricting flexibility and hindering systematic evaluation across alternative AI services \cite{sperber2020endtoendpromise}.

Motivated by these challenges, this thesis investigates the design and implementation of a real-time AI-based simultaneous interpretation system with a strong emphasis on system architecture, latency control, and practical integration. Rather than focusing on training new models, the work adopts a system-engineering perspective that leverages state-of-the-art AI services within a unified, modular pipeline.  The work is guided by three main objectives:
\begin{enumerate}
  \item To design a real-time AI-based simultaneous interpretation pipeline that processes continuous speech input using streaming ASR, NMT, and TTS synthesis, while ensuring low end-to-end latency, stable incremental output, and support for modular, provider-agnostic, and multi-language operation.
  
  \item To develop and apply a system-level evaluation framework that quantifies the performance of the proposed pipeline, with particular emphasis on end-to-end latency, component-level delays, and the impact of segmentation parameters on translation behavior.
  
  \item To investigate the feasibility of integrating AI-based interpretation pipelines with professional interpretation hardware and audio systems, considering standardized audio formats, real-time routing requirements, and practical deployment constraints.
\end{enumerate}

The remainder of this thesis is structured as follows: 
Chapter~\ref{chapter:background} provides background information and reviews related work in traditional interpretation systems, core speech and language technologies, and existing AI-based interpretation approaches, culminating in the identification of key limitations and research gaps. 
Chapter~\ref{chapter:methodology} presents the methodology and system design, detailing the architecture of the proposed real-time interpretation pipeline, its individual components, and the evaluation framework used to assess system performance. 
Chapter~\ref{chapter:results} presents the experimental results, focusing on end-to-end latency behavior, component-level delays, and system performance under real-time conditions, while Chapter~\ref{chapter:discussion} provides a detailed discussion and interpretation of these findings.
Chapter~\ref{chapter:hardware} examines the feasibility of integrating AI-based interpretation pipelines with professional interpretation hardware and audio systems, addressing practical deployment considerations and interface requirements. 
Finally, Chapter~\ref{chapter:conclusion} concludes the thesis by summarizing the main findings and outlining directions for future work.
